\section{Convex product quantizers and the exterior carrier}\label{sec:quantizers}

Throughout this section $p>2$.  We use the carrier
$F_b$, its exact target set $O$, the smaller exact set $O_1$, and the
protected endpoint boxes supplied by Lemma~\ref{hp:wells} and
Proposition~\ref{hp:carrier}.  We construct a Lipschitz target map $T$
whose three-dimensional input
blocks lie in $O$, where the carrier is already exact.  Outside those
blocks, $T$ takes values in a polyhedral two-complex.  The retained
deficient-rank gradients must survive on its faces.  Small central
prisms will then be removed from the input space to produce an exterior
with standard annuli and disks.  A final local replacement makes the
carrier exact near this new boundary, ready for the source-wall
construction.

The carrier $F_b$ is a proper continuous surjection in
$W^{1,p}_{\mathrm{loc}}$ and has the relative opening property of
Lemma~\ref{hp:opening}.  On the inverse of $O$ it agrees locally with
specified affine or PL homeomorphisms.  The compact retained rank data
are denoted by $K$; near $K$ the carrier equals $F_\kappa$.  All
smallness prescriptions may vary continuously toward the boundary.
A locally finite prescription therefore imposes only finitely many
accuracy requirements on each compact set.

The construction combines classical convex conjugacy and proximity
\cite{Moreau1965} with the polyhedral geometry of power diagrams
\cite{Aurenhammer1987}.  The contact exclusion, protected paired
arrays, orthogonal input products and properness estimates required
here are established below; they are not consequences of those
general constructions alone.

\subsection{Convex contacts and product cells}

At a smooth contact point, a function minus its affine support has a
local minimum, so its Laplacian is nonnegative.  We therefore force
contacts into $O$ by making the Laplacian negative outside $O$.  The
small residual volume allows this without losing the protected paired
supports or appreciably changing the potential.

\begin{lemma}[Excluding contacts by a Poisson perturbation]\label{qt:poisson}
Let $b_1$ be the protected-well perturbation of Lemma~\ref{hp:wells}.
Fix a positive $1$-Lipschitz function $\delta$ on $\Lambda$, extended by
zero outside, with $\delta\leq\dist(\cdot,\R^3\setminus\Lambda)$.
The well sizes can first be chosen sufficiently small, and the residual
volume prescriptions for $\Lambda\setminus O$ in
Proposition~\ref{hp:carrier} can subsequently be chosen sufficiently small,
so that there is $b_2\in C^\infty(\Lambda)$, zero outside $\Lambda$,
with the following properties.  For an absolute $C$,
\begin{equation}\label{qt:laplacian}
 \psi(x)=\tfrac12|x|^2+b_1(x)+b_2(x),\qquad
 \Delta\psi\leq C\text{ in }\Lambda,\qquad
 \Delta\psi<0\text{ in }\Lambda\setminus O.
\end{equation}
The function $b_2$ vanishes on neighborhoods of all protected endpoint
boxes, every prescribed affine support indexed by a protected
horizontal slope remains a global support of $\psi$, and, with
arbitrarily small $\epsilon>0$,
\begin{equation}\label{qt:potential-smallness}
 |b_1+b_2|\leq\epsilon\delta^2,\qquad
 b_1+b_2=o\bigl(\dist(x,\R^3\setminus\Lambda)^2\bigr)
 \quad(x\to\partial\Lambda).
\end{equation}
\end{lemma}

\begin{proof}
Choose closed Whitney cubes $Q$ with interior enlargements $Q^+$ and
bounded overlap of the enlargements.  For each cube the disallowed set
$E_Q=Q\setminus O$ is compact.  Its measure is smaller than a threshold
that has not yet been specified.  Smooth a neighborhood of $E_Q$ inside
$Q^+$ to obtain $0\leq\rho_Q\leq1$, equal to one near $E_Q$ and with
arbitrarily small support measure in units of $Q$.  This is possible by
outer regularity, provided the residual threshold is small enough.
In normalized coordinates put this density in a fixed three-dimensional
torus $\mathbb T$ containing $Q^+$, and solve
\[
 \Delta u_Q=-A\rho_Q+
       \frac{A}{|\mathbb T|}\int_{\mathbb T}\rho_Q,\qquad
 \frac{1}{|\mathbb T|}\int_{\mathbb T}u_Q=0.
\]
Here $A$ is fixed, larger than $3+\sup\Delta b_1+2$.
For completeness, the required smallness is already a consequence of
the periodic Green kernel: its singularities are $O(|x|^{-1})$ and those
of its first derivative are $O(|x|^{-2})$.  Both kernels belong to
$L^{q/(q-1)}$ for fixed $q>3$.  Convolution and H\"older's inequality
therefore give
\[
 \|u_Q\|_\infty+\|Du_Q\|_\infty
 \leq C_q A\|\rho_Q\|_{L^q}
 \leq C_q A|\supp\rho_Q|^{1/q}.
\]
The mean-zero Green solution is smooth because the density is smooth.
Multiply it by a cutoff equal to one on $Q$ and supported in $Q^+$,
and return to physical coordinates with the factor $\ell(Q)^2$.
The resulting $v_Q$ satisfies $\Delta v_Q\leq\varepsilon_Q$
everywhere and $\Delta v_Q\leq-A+\varepsilon_Q$ on $E_Q$;
its value and gradient are as small as desired in the corresponding
scaled units.  The errors here include the mean compensation and both
cutoff terms, which involve only $u_Q,Du_Q$.

There is a smooth cutoff $\chi$, zero on neighborhoods of the compact
protected endpoint boxes and one outside $O_1$, whose transition is
contained in $O_1$.  Its derivatives on each $Q^+$ are fixed before the
residual thresholds.  Set $b_2=\chi\sum_Qv_Q$.  The sum is locally
finite.  In
\[
 \Delta(\chi v_Q)=\chi\Delta v_Q+2D\chi\cdot Dv_Q+v_Q\Delta\chi
\]
the possibly large negative term is harmless for an upper bound, and
all other terms can be made arbitrarily small per cube.  Bounded overlap,
with the errors chosen below $1$ in total at every point, proves the
first inequality in \eqref{qt:laplacian}.  At a point outside $O$ the
cutoff is identically one locally and at least one $E_Q$ contributes
$-A$; this proves the strict negative inequality.  Make the value bounds
on each enlargement smaller also than the prescribed local multiple of
$\delta^2$, and let that multiple tend to zero with boundary distance.
This gives \eqref{qt:potential-smallness}, after the initial analogous
choice for $b_1$.

It remains to check the global supports.  For a fixed compact horizontal
endpoint box, subtract one of its prescribed affine supports from
$|x|^2/2+b_1(x)$.  Its only zeros are the two prescribed endpoint points
for that slope, by the two-well test.  Uniformly over the compact set of
slopes, it has a strictly positive minimum on every compact set outside
the endpoint neighborhoods.  At infinity the gap grows quadratically.
The added perturbation vanishes in the endpoint neighborhoods.  Choose
its absolute value elsewhere below half the preceding gap, using the
minimum over the compact slope set.  There are only finitely many protected
patches, by Lemma~\ref{hp:wells}, and their supports have been separated.
The minimum of their positive gap requirements is positive away from
the endpoint neighborhoods.  Taking the smaller of these requirements preserves all supports
simultaneously.  All these requirements determine only smaller residual
thresholds, after the endpoint neighborhoods have been fixed, as allowed
by Proposition~\ref{hp:carrier}.
\end{proof}

\begin{lemma}[Contact curvature and dual separation]\label{qt:contact}
Let $\phi$ be the lower convex envelope of $\psi$ from
Lemma~\ref{qt:poisson}, and let
\[
 F(q)=\phi^*(q)=\sup_x\{q\cdot x-\psi(x)\}.
\]
There is an absolute $C_0$ such that $\phi$ is $C^{1,1}$ and
$0\leq D^2\phi\leq C_0 I$ in the distributional sense.
With $m=1/C_0$ reduced if necessary, every actual contact maximizer
$x_i$ at slope $q_i$ satisfies
\begin{equation}\label{qt:dual-gap}
 F(q)\geq F(q_i)+x_i\cdot(q-q_i)+\tfrac m2|q-q_i|^2.
\end{equation}
For any requested fine motion tolerance the perturbations can be chosen
so that all contact maximizers at $q$ lie within that tolerance of $q$.
For $q\notin\Lambda$ the unique maximizer is $q$ and $F(q)=|q|^2/2$.
For $q\in\Lambda$ every maximizer lies in $O$; over a compact subset of
$\Lambda$ their union is compactly contained in $O$.
\end{lemma}

\begin{proof}
The function $\psi$ has quadratic growth at infinity.  A supporting
plane of its convex envelope consequently has a compact, nonempty set
of contacts, and each point of the envelope on that plane is a convex
combination of at most four contacts.  One way to see attainment is to
minimize the average of $\psi$ over probability measures of fixed
barycenter.  Quadratic growth gives tightness and a bounded second
moment; passage to the limit gives a minimizer.  Separation by a
supporting plane forces its support onto the contact set, and
Carath\'eodory's Theorem reduces it to four points.

At a contact in $\Lambda$, the smooth function $\psi$ minus its
supporting plane has a minimum.  Thus $D^2\psi$ is nonnegative there;
\eqref{qt:laplacian} bounds each eigenvalue above by $C$.  At a contact
on or outside $\partial\Lambda$, the zero extension and
\eqref{qt:potential-smallness} give the quadratic Taylor expansion of
$|x|^2/2$ with an $o(|v|^2)$ remainder.  If
$x=\sum\alpha_jx_j$ is a contact combination, use $x_j+v$ and $x_j-v$
as competing combinations to obtain
\[
 \limsup_{v\to0}
 \frac{\phi(x+v)+\phi(x-v)-2\phi(x)}{|v|^2}\leq C_0.
\]
Here $C_0$ is absolute and independent of the contact combination.
On any line this implies semiconcavity with constant $C_0$: if, after
subtracting $(C_0+\varepsilon)t^2/2$, the function dipped below one of
its chords, its chord-subtracted minimum would be interior.  The
symmetric second difference at that minimum is nonnegative, contrary
to the displayed strict negative upper bound.  Let $\varepsilon\downarrow0$.
Convexity and semiconcavity imply differentiability and the global bound
\[
 \phi(x+v)\leq\phi(x)+D\phi(x)\cdot v+\tfrac{C_0}2|v|^2.
\]
At a contact $x_i$, $D\phi(x_i)=q_i$.  Substituting
$v=(q-q_i)/C_0$ into the definition of the conjugate proves
\eqref{qt:dual-gap}.

Write $b=b_1+b_2$ and $d=|x-q|$ for a maximizing contact.  Comparison
with $x=q$ gives
\[
 \tfrac12d^2\leq |b(x)|+|b(q)|
 \leq\epsilon\{(\delta(q)+d)^2+\delta(q)^2\}.
\]
For sufficiently small $\epsilon$ this yields
$d\leq C\sqrt\epsilon\,\delta(q)$.  When $q\notin\Lambda$ it gives
$d=0$.  Choosing the positive Lipschitz minorant $\delta$ and then
$\epsilon$ proves the fine motion assertion and keeps all contacts
inside $\Lambda$ for interior $q$.  A contact in $\Lambda\setminus O$
would have nonnegative Hessian and therefore nonnegative Laplacian,
contrary to \eqref{qt:laplacian}.  Finally, maximizing contacts form a
closed graph and remain in a bounded set over bounded slopes.  For
slopes in a compact subset of $\Lambda$ their union is therefore
compact; since it is contained in $O$, its clearance from $O^c$ is
positive.
\end{proof}

\begin{proposition}[A proper product quantizer]\label{qt:quantizer}
Fix $0<c<m$ with $c<1$.  One can choose a locally finite family of
interior sites in $\Lambda$, with one contact at each ordinary site
and the two protected contacts at each paired site.  The family can
include the rectangular paired arrays constructed in
Lemma~\ref{qt:calibration}.  Index these site and contact pairs by
$i$, writing them as $(q_i,x_i)$, and define
\begin{equation}\label{qt:obstacle}
 G(q)=\sup_i\{F(q_i)+x_i\cdot(q-q_i)+\tfrac c2|q-q_i|^2\},
 \qquad T=(\partial G)^{-1},
\end{equation}
where the supremum also includes the support based at every
$q_i\notin\Lambda$.  Here $q$ is an output of $T$, and
$y\in\partial G(q)$ is an input.  Then $T:\R^3\to\R^3$ is single-valued and
$c^{-1}$-Lipschitz, equals the identity outside $\Lambda$, and restricts
to a proper onto map of $\Lambda$ to itself, as finely close to the
identity as prescribed.

Inside $\Lambda$, $G-c|q|^2/2$ is a locally finite maximum of affine
functions.  Its diagram has compact convex polyhedral cells.  Write
$p_i=x_i-cq_i$.  Over the relative interior of a face $P$, whose active
indices are $I(P)$, its exact input product is
\begin{equation}\label{qt:product}
 \partial G(q)=cq+\conv\{p_i:i\in I(P)\},\qquad
 T(cq+v)=q\quad(q\in\relint P).
\end{equation}
The two factors have orthogonal direction spaces of complementary
dimensions.  Every retained support owns a genuine three-dimensional
output cell $C_i$.  Its closed input block
\begin{equation}\label{qt:blocks}
 X_i=cC_i+p_i
\end{equation}
is compactly contained in $O$.  These closed blocks are pairwise
disjoint and have locally finite, mutually disjoint protective
neighborhoods within the exact regions of the carrier.  Outside their
union the output of $T$ lies in the two-skeleton of the diagram.
\end{proposition}

\begin{proof}
\emph{Localizing the active supports.}
For a support $L_i$ appearing in \eqref{qt:obstacle},
\eqref{qt:dual-gap} gives
\begin{equation}\label{qt:active-bound}
 L_i(q)\leq F(q)-\tfrac{m-c}{2}|q-q_i|^2.
\end{equation}
A sufficiently fine locally finite net gives $G\geq F-e$ for any
prescribed positive continuous error $e$ on $\Lambda$: for each compact
set, contacts and slopes are bounded, so the support at a nearby site
converges uniformly to $F$ there.  A Whitney refinement with compact
buffers gives these estimates simultaneously.  Also $G\leq F$.
Outside $\Lambda$ its own support is present, so $G=F=|q|^2/2$ there.
Choose $e(q)$ so small that
\[
 \tau(q)=\sqrt{2e(q)/(m-c)}<\tfrac14\dist(q,\Lambda^c),
\]
and impose any additional fine upper bounds on $\tau$ needed below.
Every active site is within $\tau(q)$ of $q$ by
\eqref{qt:active-bound}.  Exterior supports and sites outside that
compact interior neighborhood are uniformly below the supremum.
The remaining interior sites are finite locally, so the supremum is
attained and gives the asserted polyhedral diagram.  This argument
also applies on all closed cell boundaries.

\emph{The proper Lipschitz map.}
The function $G$ is $c$-strongly convex and grows quadratically.
Consequently $G(q)-y\cdot q$ has a unique minimum for every $y$.
Strong monotonicity gives
\[
 (y-y')\cdot(T(y)-T(y'))\geq c|T(y)-T(y')|^2,
\]
proving the Lipschitz bound.  At a point outside $\Lambda$, equality
$G=F$ and differentiability of $F$ imply $\partial G(q)=\{q\}$:
every supporting affine function for $G$ there also supports $F$.
Thus $T$ is the identity outside, and no interior point is sent outside.
For interior $q$, \eqref{qt:product} expresses each input as a convex
combination of
\[
 cq+p_i=x_i+c(q-q_i).
\]
The contact motion bound and the bound $|q-q_i|\leq\tau(q)$ make every
such point arbitrarily finely close to $q$; in particular it lies in
$\Lambda$.  This proves surjectivity onto $\Lambda$.  Arrange
$|T(y)-y|<\tfrac14\dist(y,\Lambda^c)$, with analogous inverse motion
control in terms of $q$.  A sequence of inputs approaching
$\partial\Lambda$ then has outputs approaching that boundary.  Since
$\Lambda$ is bounded, inverse images of compact subsets are compact,
which proves properness.

\emph{Orthogonal products and exact input blocks.}
Subtracting the common quadratic in \eqref{qt:obstacle} gives slopes
$p_i$.  Equalities between its active affine functions have normals
$p_i-p_j$.  The tangent space of a relatively open face is the common
kernel of these equalities, so their span is its entire normal space.
This proves the complementary orthogonal product assertion and the
subgradient formula.  The closure of a product over $P$ is meant here
when a face is closed; a boundary point can have further incident
products in its complete fibre.

An ordinary support at its own site is strictly above every support
from another site, by \eqref{qt:active-bound}.  The gap is uniformly
positive locally: nearby distinct sites are separated, and the
remaining sites, including the exterior family, have positive distance
from that site.  It therefore wins on an open set.  At a paired site
the two contacts have different normal components; each wins on one
open side of their common tie plane while the same other-site gap
persists.  Thus all retained supports own full cells.

Two offsets from different sites cannot coincide.  Indeed strong
monotonicity of $\partial F$, obtained by adding the two inequalities
\eqref{qt:dual-gap}, yields
\[
 (x_i-x_j)\cdot(q_i-q_j)\geq m|q_i-q_j|^2.
\]
If $p_i=p_j$, its left side is $c|q_i-q_j|^2$, a contradiction.
At one site discard identical contacts.  If two closed blocks met,
uniqueness of $T$ would give the same output $q$ and then the same
offset, another contradiction.

Choose $\tau$ smaller also than the local contact clearance in $O$
divided by $2c$, using compact neighborhoods of the slopes.
Then every $x_i+c(q-q_i)$ for $q\in C_i$ lies in $O$.  A fixed cell
cannot approach $\partial\Lambda$, since
$|q-q_i|<\dist(q,\Lambda^c)/4$ throughout it.  Boundedness makes $C_i$
and $X_i$ compact.  Over a compact input set, motion control confines
outputs and active sites to compact interior sets, so only finitely
many blocks occur.  Pairwise disjoint compact blocks in this locally
finite family admit pairwise disjoint neighborhoods in $O$, reduced
further to the prescribed exact carrier neighborhoods.  Finally, over
the interior of a full output cell the product is the singleton graph
$y=cq+p_i$, which is precisely its full input block.  This proves the
last assertion.
\end{proof}

\subsection{Calibration and strict approximation}

For the rest of the construction fix $c$ universally, for example
$c=\min\{m,1\}/2$.  It is independent of every later accuracy
parameter; constants depending only on this choice are absolute.

\begin{lemma}[Paired arrays, redistribution, and phase calibration]\label{qt:calibration}
The choices in Proposition~\ref{qt:quantizer} can be made together
with a locally bi-Lipschitz homeomorphism
$J:\Lambda\to\Lambda$, equal to the identity near every closed $X_i$,
such that $TJ$ has an absolute Lipschitz bound.  Put
\[
 \widehat F=JF_b,\qquad q_0=T\widehat F.
\]
Apart from arbitrarily small prescribed source weight, $q_0$ maps the
retained rank-one and rank-two data into the interiors of horizontal
rectangular true two-faces.
The carrier $\widehat F$ retains the exact data over the blocks and its
relative opening property, and $|Dq_0|\leq C|DF_b|$ almost everywhere.
On the retained compact rank-$k$ data, $k=1,2$,
\begin{equation}\label{qt:calibrated-gradient}
 Dq_0=a_*^{-1}\operatorname{proj}_{n^\perp}DF_\kappa,
 \qquad |Dq_0-Df|\leq C(\lambda+1-a_*)|Df|,
\end{equation}
where $c<a_*<1$ may be as close to one as required and $\lambda$ is
the preceding flattening accuracy.  The rank is exactly $k$ and
$\ker Df\subset\ker Dq_0$.  For rank one the rectangle can have any
previously prescribed small short-to-long aspect, with its first axis
arbitrarily close to the range line of $Dq_0$.  Further compact margins,
sector avoidance, and uniform jet tests may be imposed by trimming.
\end{lemma}

\begin{proof}
On a protected patch write $q=z+h+q_v n$.  At the sites $q_i=z+u_i$ of
a Cartesian lattice in $n^\perp$, with periods $d_1,d_2$, retain both
contacts $x_i^\pm=z+u_i\pm d_\perp n$.  The protected support identity is
\[
 F(z+u)=\tfrac12|z+u|^2-d_\perp^2.
\]
Exclude nonmatching sites from a narrow tube about a compact interior
part of the plane $q_v=0$.  The tube width is chosen after the error and
active-site tolerances in the preceding proof, and the periods are
chosen much smaller still.  A point in the tube has a paired site at
distance at most its width plus $\sqrt{d_1^2+d_2^2}$; local continuity
and the support formula give the required dual approximation there.
Outside use a sufficiently fine allowed net.  Thus exclusion does not
invalidate the previous approximation.

More explicitly, on the central plane the deficit of the nearest pair
is $O(d_1^2+d_2^2)$, whereas every nonmatching site has, on a smaller
protected interior, a fixed positive distance and hence a fixed
positive deficit by \eqref{qt:active-bound}.  Decrease the periods until
the former is smaller.  Dominance persists in a vertical neighborhood.
The affine parts of two matching pairs differ there by the nearest-node
quadratic comparison with coefficient $1-c>0$.  Their bisectors are
therefore exactly the Cartesian bisectors, and the two signs tie
exactly on $q_v=0$.

The subgradient formula now gives, for protected inputs
$y=z+h+vn$ with $|v|<d_\perp$,
\begin{equation}\label{qt:scalar-quantizer}
 T(y)=z+(Q_{c,1}(h_1),Q_{c,2}(h_2))_{\rm horiz}.
\end{equation}
In a period centered at node $u$ of pitch $d$, the scalar map is
$u+(h-u)/c$ for $|h-u|\leq cd/2$, and is constant, equal to the
appropriate output period endpoint, on each remaining half-gap.
At a tie all tied nodes have both signs, so their convex hull includes
the entire interval $|v|\leq d_\perp$; thus the formula also holds on
the seams.  The pertinent full output rectangles in $q_v=0$ are true
faces of the diagram, not arbitrary two-dimensional subsets of cells.

The scalar map $Q_c$ has slope $1/c$ on only a fraction $c$ of each
input period and collapses the complementary gaps.  We enlarge this
fraction to $a$ close to one, so that most retained data see the nearly
unit slope $1/a$.  For this purpose choose a Lipschitz cutoff
$a(h,v)\in[c,a_*]$, equal to $a_*$ on a
neighborhood of the protected values and equal to $c$ off a compact
subset of the separability region $|v|<d_\perp$.  In a period define
$R_a$ to map the interval of fraction $a$ linearly onto the interval of
fraction $c$, and the two complementary intervals linearly onto their
counterparts, fixing the endpoints.  Thus $Q_cR_a=Q_a$, where $Q_a$
is the same scalar map with $c$ replaced by $a$.
We must also check that the two simultaneous redistributions remain
invertible when $a$ varies with the input.  In centered coordinates
the inverse of $R_a$ is
\[
 B_a(w)=
 \begin{cases}
 (a/c)w,&|w|\leq cd/2,\\
 ad/2+\dfrac{1-a}{1-c}(w-cd/2),&cd/2\leq w\leq d/2,
 \end{cases}
\]
with the odd extension on the negative half-period.  Hence
$|\partial_a B_a|\leq d/2$, including across the fixed inverse
breakpoints.  Define simultaneously
\[
 J(h,v)=\bigl(R_{a(h,v),1}(h_1),R_{a(h,v),2}(h_2),v\bigr).
\]
For given output $(w,v)$, its inverse solves
\[
 h_j=B_{a(h,v),j}(w_j),\qquad j=1,2.
\]
Choosing the periods so that
$\sqrt{d_1^2+d_2^2}\Lip(a)<1/2$ makes this a contraction on $\R^2$,
after extending the periodic formulas and putting $a=c$ outside the
support.  It has a unique solution, Lipschitz in $w,v$.  The forward
formulas are Lipschitz for the fixed $a_*<1$; consequently $J$ is
bi-Lipschitz on this patch.  It is the identity off the patch and maps
it onto itself.  The disjoint patches assemble locally finitely.

The composite $Q_cR_a=Q_a$ has ordinary input slopes bounded
by $1/c$ and parameter derivative bounded by $C_cd$; the latter follows
either from its formula $u+(h-u)/a$ on the middle interval or from
continuity at the moving breakpoints.  The chosen period bound makes
$TJ$ uniformly Lipschitz, independently of how close $a_*$ is to one.
The support is strictly between the two endpoint heights; a full block
there would require an extreme sign, at height $\pm d_\perp$.
Thus the support avoids all closed full blocks, with a neighborhood.
The same is true of its image.  The carrier opening property transfers
under the target homeomorphism $J$, and the derivative bound follows
from the Lipschitz chain rule.

All these choices can be uniform over lattice phases.  Fix nested
horizontal boxes
$H_{\rm inner}\Subset H_{\rm array}\Subset H_{\rm endpoint}$,
the excluded tube, and the cutoff support first.  For every phase
retain all lattice nodes in $H_{\rm array}$, taking both periods below
one tenth of the nesting margins.  The nearest-node deficit is bounded
by $C(d_1^2+d_2^2)$ for every phase, whereas the nonmatching-site gap
on the inner box is fixed.  Thus a common separability slab
$|v|<d_\perp-\sigma$ contains the fixed cutoff support for all phases.
Its separation from full blocks and the inverse contraction estimate
are phase-independent.

Choose the phases after the compact data, cutoffs, and periods have
been fixed.  For any fixed datum, uniform independent translations of
the two lattice coordinates place it in their middle intervals with
probability $a_*^2$.  Fubini's Theorem for the finite source weight
therefore gives a phase losing at most the prescribed multiple of
$1-a_*^2$, with room for further losses.  Conditional on membership
in these intervals, its two normalized output coordinates are uniform
on the output rectangle: this follows directly from the linear
formula and translation-invariance of phase measure.  Consequently
any prescribed zero-area subset of a normalized rectangle, such as
its center, perimeter, or finitely many sector rays, is avoided for
almost every retained datum for almost every phase.  Neighborhoods
of those sets have losses tending to zero by the same averaging.
This argument uses no absolute continuity of the projected data
measure.  Independently shifted finer rectangular subdivisions give
the same conclusion after discarding fitting margins.

Where the retained data have $a=a_*$ and both middle intervals are
strict, differentiation of \eqref{qt:scalar-quantizer} gives
\eqref{qt:calibrated-gradient}; the preceding flattening bound and
$|n-e|\leq C\lambda$ give its error estimate.  The Lipschitz
composition annihilates every direction in $\ker Df$ at a source
point where the compared jets exist.  After the preliminary positive
singular-value bounds have been fixed, choose the errors smaller than
those bounds, so the rank cannot decrease.  It cannot increase because
of the kernel inclusion.  The first lattice axis follows the selected
rank-one range direction; its angular errors were chosen smaller than
the assigned aspect.  Finally continuity, regularity of the finite
source weight, and differentiation allow compact trimming to all
claimed interior and uniform jet margins.
\end{proof}

\begin{lemma}[Strict quantizers]\label{qt:strict}
The map $T$ has proper locally bi-Lipschitz approximations
$T_\eta:\Lambda\to\Lambda$ with the same target.  Given any positive
continuous local tolerance $\epsilon_0(y)$, they can be chosen with
\[
 |T_\eta(y)-T(y)|<\epsilon_0(y),\qquad
 |DT_\eta(y)-DT(y)|<\epsilon_0(y)\quad\text{for a.e. }y.
\]
On every product piece the derivatives of $T$ are its affine
tangential derivatives, agreeing on directions tangent to common
strata.  If $h$ is locally bi-Lipschitz, the pushforward by $Th$ of
volume restricted to each product interior is absolutely continuous
with respect to the dimensional measure on its output stratum.
\end{lemma}

\begin{proof}
For $q=T(y)$ set
\[
 T_\eta(y)=q+\eta(q)(y-q),
\]
where $\eta>0$ is smooth, bounded above by a small constant, and has a
small global Lipschitz constant.  Such functions with arbitrary local
upper bounds on their value and first derivative are obtained from a
locally finite smooth partition of unity: choose the coefficient of
each term smaller than all required bounds on its support, divided by
the local overlap count and the first derivative bound of that term.
Put $M=\sup_{\Lambda}|y-T(y)|<\infty$.  For $q_j=T(y_j)$,
monotonicity and direct subtraction give
\[
 (T_\eta(y_1)-T_\eta(y_2))\cdot(q_1-q_2)
 \geq(1-\|\eta\|_\infty-M\Lip\eta)|q_1-q_2|^2
 \geq\tfrac12|q_1-q_2|^2.
\]
Thus differences of outputs control differences of $q$.  Subtracting
again in the defining formula, and using a positive lower bound for
$\eta$ on each compact set, controls $|y_1-y_2|$ locally by the output
difference.  This proves injectivity and local inverse Lipschitz
bounds, including the case $q_1=q_2$, when the output difference is
exactly $\eta(q_1)(y_1-y_2)$.  Forward Lipschitz bounds follow from those
of $T$ and $\eta$.  Fine motion keeps the segment from $y$ to $q$
inside $\Lambda$ and makes $T_\eta$ proper, as in
Proposition~\ref{qt:quantizer}.  Invariance of domain gives an open
image; properness gives a relatively closed image.  Connectedness of
$\Lambda$ proves surjectivity.

At almost every input,
\[
 D(T_\eta-T)=\eta(T)(I-DT)
 +(y-T)\otimes\bigl(D\eta(T)DT\bigr).
\]
Properness first transfers any local input tolerance to a local output
tolerance; choose $\eta,D\eta$ below it in the preceding construction.
This proves both requested estimates.  On a product piece
\eqref{qt:product} is orthogonal projection to its face, with factor
$c^{-1}$ in tangential directions.  Its restriction to a common stratum
is the same map from either side, proving tangential agreement.
Fubini on the product, followed by the change-of-variables inequality
for a locally bi-Lipschitz $h$, proves the dimensional absolute
continuity assertion.  At a zero-dimensional output stratum the map
$Th$ is constant and its weak derivative vanishes almost everywhere
on that level set.
\end{proof}

\subsection{Face coordinates and the graph exterior}

\begin{lemma}[Subdividing the true faces]\label{qt:subdivision}
The true two-faces of the quantizer diagram have locally finite,
conforming subdivisions into convex polygonal subcells $D$ with
chosen centers $d_D$.  Protected rectangular faces may be kept as
rectangles or subdivided by finer rectangular grids.  Elsewhere the
ratio of diameter to center-boundary clearance has an absolute bound
away from arbitrarily small prescribed neighborhoods of the original
vertices.  Inside those neighborhoods it has a finite bound depending
only on the original face.  The latter constants can be confined to
regions of arbitrarily small weighted $\int|Dq_0|^p$ cost.  Additional
rays from a protected center to inserted perimeter vertices can be
chosen to avoid almost every retained datum.
\end{lemma}

\begin{proof}
In a compact convex face $P$, choose a maximal separated net of spacing
$d$ and intersect its planar Voronoi cells with $P$.  Maximality bounds
each cell's diameter by $C d$, while separation gives the intersection
of a disk of radius $c d$ about its generator with $P$.  Fix an interior
disk of $P$ once.  Interpolating the generator toward its center by a
distance proportional to $d$ gives a disk of radius $c_Pd$ in that
intersection: convexity carries the fixed interior disk into $P$, and
its interpolated center and radius both remain within the small disk
about the generator.  Thus the clearance ratio is bounded by a finite
constant depending on $P$.  Outside a prescribed neighborhood of the
vertices and for sufficiently small $d$, at most one supporting edge
line cuts this small disk.  A disk cut by one half-plane through or
outside its center contains a disk of one fixed fraction of its radius.
The clearance constant there is absolute.  Choose the center $d_D$
of each cell in the indicated inscribed disk.

Match the subdivisions of a shared original edge by overlaying the
finitely many incident subdivisions.  This inserts only collinear
perimeter vertices and does not change any cell or its clearance.
Choose the finite nets generically first, so no Voronoi vertex lies
on an original edge and each incident Voronoi segment meets that edge
transversely.  The finitely many nonzero angles and strict convexity
margins have a positive minimum.  Changes of division positions on
original edges and the corresponding changes in adjacent segments
are taken below these margins and preserve cyclic order.  They
therefore preserve convexity and the clearance bounds.  Collinear
overlay nodes may slide along a side without changing its polygon.  For a fixed interior point of a
protected rectangle, the ray from its center through that point
intersects the perimeter at one location.  A continuously distributed
perturbation of an extra perimeter vertex equals that location with
probability zero.  Fubini for the finite retained measure gives
simultaneous avoidance, even for singular data.  The original center,
corner, and sector-ray exclusions have already been arranged by
Lemma~\ref{qt:calibration}.  All choices are finite per face and locally
finite in $\Lambda$.

A Sobolev map has zero weak derivative almost everywhere on any one
of its level sets: apply the scalar statement to each coordinate,
which follows by truncating that coordinate to a shrinking interval
about the level and using the Sobolev chain rule.  In particular
$Dq_0=0$ almost everywhere on $\{q_0=v\}$ for an original vertex $v$.
On a compact localization, dominated convergence gives
\[
 \lim_{\rho\downarrow0}
 \int_{\{|q_0-v|<\rho\}}|Dq_0|^p=0.
\]
The map $q_0$ is proper and locally Sobolev, so the localizations needed
for a compact target neighborhood are compact source localizations.
Fix the original-face clearance factors first, and then choose the
vertex neighborhoods to make these integrals smaller than any given
error divided by those factors.  This also permits any previously
fixed collar weights.  Use a summable family of errors for the locally
finite vertices.  Only after these neighborhoods have been fixed take
$d$ sufficiently small.  The arbitrary original-face constants thus
multiply only the separately paid vertex contributions.
\end{proof}

\begin{lemma}[Polar coordinates and their sharp rectangular bound]\label{qt:polar}
For each subcell $D$, put
\[
 R_{D,e}(t,s)=d_D+r\bigl(\xi_e(s)-d_D\bigr),\qquad
 r=\exp(t/L_D),
\]
where $\xi_e$ is arclength on a perimeter subedge and the same arclength
label is used in all incident sectors.  Choose $L_D$ comparable to the
maximum radius of $D$; in a protected rank-one rectangle choose it
as the long-axis half-length and choose $d_D$ to be the geometric
midpoint of the rectangle.  On an open sector,
\begin{equation}\label{qt:polar-bounds}
 |D_qt|+|D_qs|\leq C_D/r,\qquad
 |\partial_tR_{D,e}|+|\partial_sR_{D,e}|\leq C r.
\end{equation}
The first constant depends only on the clearance ratio, and the second
is absolute.  For a rectangle of half-length $L$ and half-width
$\gamma L$, and its unit long-axis vector $v_1$,
\begin{equation}\label{qt:rank-one-polar}
 |D_qt(v_1)|+|D_qs(v_1)|\leq(1+\gamma)/r,
 \qquad
 \max\{|\partial_tR|,|\partial_sR|\}
 \leq\sqrt{1+\gamma^2}\,r.
\end{equation}
\end{lemma}

\begin{proof}
Translate $d_D$ to zero.  If the supporting line of $e$ is
$\nu\cdot\xi=b>0$, then $r=(\nu\cdot q)/b$,
$Dr=\nu/b$, and $Dt=L_D\nu/(br)$.  Also
\[
 \xi'(s)D s=\frac1r\left(I-\frac{q\otimes\nu}{br}\right),
\]
when restricted to the sector plane.  Since $b$ is bounded below by
the center clearance and $|q|/r$ by the maximum radius, these are the
first estimates.  Directly,
$\partial_tR=r\xi/L_D$ and $\partial_sR=r\xi'$, proving the second.
For the rectangle, on a short end $q_1=\pm rL$, so
$|D t(v_1)|=1/r$ and $|D s(v_1)|\leq\gamma/r$.
On a long side $r$ depends only on the short coordinate, so the two
values are $0$ and $1/r$.  Finally $|\xi|\leq L\sqrt{1+\gamma^2}$
and $|\xi'|=1$.  Splitting a straight side only changes the additive
origin of $s$, so none of the estimates changes.
\end{proof}

\begin{proposition}[The graph exterior and its two colours]\label{qt:exterior}
Choose a small $r_{u,D}>0$ in every subcell.  In input target
coordinates remove the full closed blocks $X_i$ and the open cores
of the interval prisms over
$d_D+r_{u,D}(D-d_D)$, with their ends attached to the full blocks.
The closure $M_y$ of the remainder is a locally finite PL
three-manifold with boundary.  The radial projection in each subcell,
identity on subedges, defines a continuous map
\[
 \pi T:M_y\longrightarrow\Gamma,
\]
where $\Gamma$ is the subedge graph.  Each component of $M_y$ is a
possibly noncompact graph handlebody.  In particular every PL embedded compact
interior two-sphere bounds a PL ball and
\begin{equation}\label{qt:exterior-homology}
 H_2(M_y;\mathbb Z)=0.
\end{equation}
For intermediate levels there are two standard families of proper
walls: $A_D(t)$ is the perimeter at common radial label $t$, times
the face-offset interval; $B_e(s)$ is the entire inverse of the
intermediate subedge point under $\pi T$.  Walls within one colour are
disjoint.  Each $A$ is an incompressible annulus and each $B$ is a disk.
An incident pair has two boundary intersection points, one on each
annular end; a nonincident pair has none.  Consequently transverse
individually standard walls with this boundary incidence have one
spanning intersection arc for each incident pair, and otherwise only
circles; every such circle is contractible in both walls.
\end{proposition}

Figure~\ref{qt:fig-product-exterior} shows one face product and the
two wall directions.  The proof also accounts for the offset polygons
and polytopes over true edges and vertices.

\begin{figure}[tb]
\centering
\begin{tikzpicture}[font=\small,>=Latex]
\begin{scope}[x=1.3cm,y=1.3cm]
 \node[font=\bfseries] at (1.15,2.97) {Output face $D$};
 \fill[gray!5] (0,0) rectangle (2.3,1.9);
 \draw[thick] (0,0) rectangle (2.3,1.9);
 \draw[blue!70!black,thick] (.46,.38) rectangle (1.84,1.52);
 \fill[white] (.92,.76) rectangle (1.38,1.14);
 \draw[densely dashed,gray] (.92,.76) rectangle (1.38,1.14);
 \draw[red!65!black,thick] (1.38,.95)--(2.3,.95);
 \fill (1.15,.95) circle (1.1pt);
 \node[below left] at (1.15,.95) {$d_D$};
 \fill[blue!70!black] (1.84,.95) circle (1.4pt);
 \node[below left,inner sep=2pt] at (1.84,.95) {$q$};
 \fill (2.3,.95) circle (1.1pt);
 \node[right] at (2.3,.95) {$\xi_e(s)$};
 \draw[->,thick] (2.46,1.25)--(2.46,1.72);
 \node[right] at (2.46,1.52) {$s$};
 \node[blue!70!black,above] at (1.15,1.54) {$t=L_D\log r$};
 \node[align=center] at (1.15,-.45)
 {$q=d_D+r(\xi_e(s)-d_D)$};
\end{scope}
\draw[->,thick] (6cm,1.65cm)--node[above] {$T$} (4.8cm,1.65cm);
\begin{scope}[shift={(7.1cm,.35cm)},
 x={(1.3cm,0cm)},y={(.42cm,.36cm)},z={(0cm,1.28cm)}]
 \node[font=\bfseries] at (1,.7,2.55) {Input face prism};
 \fill[gray!18] (0,0,0)--(2,0,0)--(2,1.4,0)--(0,1.4,0)--cycle;
 \fill[gray!12] (0,0,1.7)--(2,0,1.7)--(2,1.4,1.7)--(0,1.4,1.7)--cycle;
 \draw[gray] (0,0,0)--(2,0,0)--(2,1.4,0)--(0,1.4,0)--cycle;
 \draw[gray] (0,0,1.7)--(2,0,1.7)--(2,1.4,1.7)--(0,1.4,1.7)--cycle;
 \draw[gray] (0,0,0)--(0,0,1.7) (2,0,0)--(2,0,1.7)
  (2,1.4,0)--(2,1.4,1.7) (0,1.4,0)--(0,1.4,1.7);
 \draw[blue!70!black,thick] (.4,.28,0)--(1.6,.28,0)
  --(1.6,1.12,0)--(.4,1.12,0)--cycle;
 \draw[blue!70!black,thick] (.4,.28,1.7)--(1.6,.28,1.7)
  --(1.6,1.12,1.7)--(.4,1.12,1.7)--cycle;
 \draw[blue!70!black] (.4,.28,0)--(.4,.28,1.7)
  (1.6,.28,0)--(1.6,.28,1.7);
 \draw[blue!70!black,densely dashed] (.4,1.12,0)--(.4,1.12,1.7)
  (1.6,1.12,0)--(1.6,1.12,1.7);
 \fill[white] (.8,.56,0)--(1.2,.56,0)--(1.2,.56,1.7)
  --(.8,.56,1.7)--cycle;
 \fill[white] (1.2,.56,0)--(1.2,.84,0)--(1.2,.84,1.7)
  --(1.2,.56,1.7)--cycle;
 \draw[gray,densely dashed] (.8,.56,0)--(1.2,.56,0)
  --(1.2,.84,0)--(.8,.84,0)--cycle;
 \draw[gray,densely dashed] (.8,.56,1.7)--(1.2,.56,1.7)
  --(1.2,.84,1.7)--(.8,.84,1.7)--cycle;
 \draw[gray,densely dashed] (.8,.56,0)--(.8,.56,1.7)
  (1.2,.56,0)--(1.2,.56,1.7)
  (1.2,.84,0)--(1.2,.84,1.7);
 \fill[red!45,opacity=.3] (1.2,.7,0)--(2,.7,0)
  --(2,.7,1.7)--(1.2,.7,1.7)--cycle;
 \draw[red!65!black,thick] (1.2,.7,0)--(2,.7,0)
  --(2,.7,1.7)--(1.2,.7,1.7)--cycle;
 \draw[blue!70!black,very thick] (1.6,.7,0)--(1.6,.7,1.7);
 \fill[blue!70!black] (1.6,.7,0) circle (1.5pt);
 \fill[blue!70!black] (1.6,.7,1.7) circle (1.5pt);
 \node[left,blue!70!black] at (-.15,0,.72) {$A_D(t)$};
 \node[right,red!65!black,align=left] at (2.12,1.4,.65)
 {$B_e(s)$\\portion};
 \node[above] at (1,.7,1.92) {end on $X_+$};
 \node[below] at (1,.7,-.22) {end on $X_-$};
 \node[rotate=90,fill=white,inner sep=1pt] at (1,.64,.85) {removed core};
\end{scope}
\end{tikzpicture}
\caption{A rectangular face and its orthogonal input product,
schematically and at independent scales; $X_\pm$ denote the adjacent
full blocks.  The quantizer $T$ collapses
the normal interval and rescales tangential directions by $c^{-1}$.
After removing the central core, a radial perimeter times the interval
is the blue annulus $A_D(t)$.  The red rectangle is only the portion
of $B_e(s)$ in this prism.  Over a true edge, the complete meridian
disk also contains its offset polygon and the spoke rectangles of all
incident faces.  The blue intersection interval joins the two ends
of the annulus on the adjacent full blocks.}
\label{qt:fig-product-exterior}
\end{figure}

\begin{proof}
Over an open true face the product formula is a polygon times an
interval perpendicular to it, whose two ends are faces of the two
adjacent full blocks.  Over an open true edge it is an interval times
a convex polygon; over a true vertex it is a convex three-dimensional
offset polytope.  These products glue by their common faces because
they are subdifferentials of the same piecewise-quadratic function.
The offset dimensions are exactly complementary by
Proposition~\ref{qt:quantizer}.  Removing the central subpolygons and
the full blocks therefore leaves the usual half-space or corner
neighborhood at each boundary point; subdivisions turn all such
corners into PL half-ball charts.  This proves the manifold assertion.
The prisms are disjoint and attach along disjoint disks after the
$r_{u,D}$ have been chosen small.  Their boundaries and the full-block
boundaries give its locally finite boundary.  The radial projection
agrees on every shared subedge.  On a true edge the map $T$ is the
projection along its polygonal offset fibre, so it gives the same
value for all incident face products.  This proves continuity of
$\pi T$ through the true edges and vertices.

Cut $M_y$ at one standard meridian $B$ over an interior point of each
subedge.  We verify that every resulting chamber is a ball.  Fix a
vertex $w$ of the subdivided graph and a true face $P$ incident there.
Inside $P$, the annular portion projecting to the truncated star at
$w$ is a disk $A_{P,w}$: for each incident subcell take the radial strip
from the first cut on one side of $w$, around $w$, to the adjacent
cut.  Successive strips meet along one radial side.  Their boundary
is a single polygonal curve alternating cut spokes and inner arcs;
if $w\in\partial P$, add the one outer boundary arc through $w$.
This is a disk fan, with empty intersection with $\partial P$ in the
interior case and one boundary arc in the boundary case.
Its product with the face-offset interval is a ball.

If $w$ is interior to a true face this product already is the chamber.
If it is interior to a true edge, start with the polygon-offset prism
over the truncated edge segment.  Each incident face product attaches
along one side disk, with disjoint relative interiors of the attaching
disks.  Attaching a ball to a ball along a boundary disk yields a ball:
identify the disk with the flat disk in a half-ball chart and flatten
the two collars, then extend the resulting boundary identification
radially.  This argument is PL after subdivision.  Thus the edge
chamber is a ball.

At a true vertex start with the full convex offset polytope.  The
incident edge products attach on its facets.  The incident fan is the
normal fan of that polytope: the equalities between active affine
supports identify exactly the facet-edge incidences in question.
After the edge attachments, each $A_{P,w}$ product attaches along the
disk over its two-sided outer boundary arc and the face-offset
interval.  The two sides meet along the corresponding polytope edge.
These are boundary disks whose relative interiors are disjoint from
the other attaching disks.  The same ball-gluing argument completes
the vertex chamber.  It also shows that its frontier meets adjacent
chambers exactly in the chosen meridian disks.  This local verification
covers degeneracies of the original diagram; it uses the full active
offset polytope rather than a simplicity assumption on the vertex.

Restoring the meridian collars therefore attaches ordinary
one-handles to disjoint zero-handles.  This is a graph-handlebody
presentation, locally finite on compact sets.  It retracts onto a
one-dimensional graph, proving \eqref{qt:exterior-homology}.
For irreducibility, a compact PL embedded sphere meets only finitely many handles;
include them and their endpoints in a finite graph chunk, cutting its
remaining handles farther away from the sphere.  Each component of
this finite chunk is an ordinary compact handlebody.  To recall the
standard elementary argument, cut such a handlebody along its
meridian disks.  Put the PL sphere transverse to them and use an innermost
intersection circle and the adjacent disk in the cut ball to remove
intersection circles by disk surgery and isotopy.  Induction reduces
to a sphere in a ball, which bounds a ball by the PL Schoenflies
Theorem; restoring each removed disk collar restores a ball on the
same side of the sphere.  Equivalently this is the induction that
attaching a boundary one-handle to a union of balls preserves
irreducibility.  The bounded ball lies in the finite chunk, hence in
$M_y$.

The $A$ product is visibly an annulus.  Its core maps under $\pi T$
to the reduced perimeter loop of $D$ in $\Gamma$.  That loop traverses
successive distinct perimeter edges without cancellation, so every
nonzero power is nontrivial in the free fundamental group of the
graph.  The induced map on the annular fundamental group is therefore
injective; in particular the annulus is incompressible.
At an artificial edge, $B$ is the union of two spoke rectangles along
their common interval, a disk.  At a true edge it consists of the
convex offset polygon and one spoke rectangle on each corresponding
boundary interval.  Gluing the rectangles to this polygon successively
along boundary intervals again gives a disk.  Their descriptions
prove disjointness in each colour and the stated boundary incidence.

For transverse individually standard walls the intersection is a
compact one-manifold, with boundary precisely the prescribed wall
boundary intersections.  An incident pair therefore has exactly one
arc, joining those two endpoints, and any remaining components are
circles.  The arc joins the two annular ends and is spanning.  Every
circle bounds a disk in the $B$ wall, hence is nullhomotopic in $M_y$;
injectivity of the annular fundamental group makes it nullhomotopic
in $A$ as well.  On a disk or annulus such a simple nullhomotopic
circle bounds a disk.  This proves the last assertions.
\end{proof}

\subsection{Exact boundary data and edge matching}

The target exterior $M_y$ now has standard annuli and disks, but its
new central-prism boundaries need not lie in the exact region of
$\widehat F$.  To pull these walls back with fixed PL boundary data,
we first make the carrier exact near those boundaries.  The changes
are confined below a prescribed radial level $a_D$, leaving the
calibrated data at higher radii unchanged.  Their new local energy
need only be finite: the radial margin lets the later target map be
constant on the changed region.  The construction must also retain enough control of
fibres to preserve sampled edge intersections when it is opened.

\begin{lemma}[A carrier exact on the central tubes]\label{qt:central-carrier}
Let $p>2$.  Let $T:\Lambda\to\Lambda$ have the properties of
Proposition~\ref{qt:quantizer}, and let
$\widehat F:\Omega\to\Lambda$ have the carrier properties retained in
Lemma~\ref{qt:calibration}.  In particular, $\widehat F$ is a proper
continuous $W^{1,p}_{\mathrm{loc}}$ carrier with connected fibres and a
countable set $E_0$ of nonsingleton values, it has the relative opening
property of Lemma~\ref{hp:opening}, and its prescribed maps over
neighborhoods of the closed full blocks $X_i$ are PL homeomorphisms.
Use the locally finite face subdivision and central prisms of
Proposition~\ref{qt:exterior}.  Fix, separately from the central radii, a
positive low activation radius $a_D$ in every subcell $D$.

The radii $0<r_{u,D}\ll a_D$ can be chosen so that the resulting closed
exterior $M_y$ has the following additional properties.  There exist a
PL homeomorphism $h_*:\Omega\to\Lambda$, a proper continuous
$W^{1,p}_{\mathrm{loc}}$ map $\overline K:\Lambda\to\Lambda$, and
\[
 F_*=\overline K h_*,\qquad M_x=h_*^{-1}(M_y),\qquad q_*=TF_* ,
\]
such that:
\begin{enumerate}
\item $\overline K$ is the identity on an open neighborhood of the
full blocks and of all the removed central prisms.  Consequently
$F_*^{-1}(M_y)=M_x$, and $F_*=h_*$ on an open neighborhood of
$\partial M_x$.  These boundary data are PL and have the originally
prescribed values near the full blocks.
\item $\overline K$, and hence $F_*$, is a fine limit of actual
homeomorphisms with these exact data.  Its fibres are nonempty and
connected, and only countably many of its target values have
nonsingleton fibres.
\item All changes from $\widehat F$ occur in compactly localized,
locally finite face neighborhoods whose old and new labels lie in
$r<a_D$.  The changes meet any prescribed fine value tolerance.
Thus $F_*=\widehat F$ in the complementary high region.  On every
compact portion of $M_x$ the labels $q_*$, and the radial and perimeter
coordinates used there, have finite ambient local Sobolev budgets.
No uniform bound for these new low-region budgets is asserted.
\end{enumerate}
There is a locally finite closed family $S$ of possible limiting inner
cube seams in the changed regions.  Away from $S$, the construction
has the local fibre formula proved below.  The portions of $S$ at which
the construction is not already exact lie away from $\partial M_x$.
\end{lemma}

\begin{proof}
\emph{Fixing the target geometry and the reference homeomorphism.}
We first fix the target geometry before choosing the final PL
approximation.  Choose slightly enlarged, pairwise disjoint polyhedral
balls around the $X_i$, contained with room in the exact PL regions of
$\widehat F$, and fix smaller open protection neighborhoods there.
These protections, the central prisms, and all boxes and sampling
margins specified next depend only on the prescribed target geometry
and on the geometric constants of Lemma~\ref{hp:cutoff}; they do not
depend on a map $h_*$ or $K_0$.

Include the properness requirements at this stage.  Fix a compact
exhaustion $C_m\subset\operatorname{int}C_{m+1}$ of $\Lambda$.
Choose a positive $1$-Lipschitz function $\tau$ so small that
\[
 \tau(z)<\tfrac14\dist(z,C_m)
 \quad\text{whenever }z\notin\operatorname{int}C_{m+1}.
\]
Such a minorant exists: only finitely many of these conditions apply
near any fixed point, and the compact gaps make each applicable
bound positive.  All geometry and approximation choices below include
the requirement that the resulting target-coordinate maps $g$, and
also $K_0$, move each $y$ by less than $\tau(y)/10$.

Over a true face, the quantizer has an orthogonal product description
with an interval as normal fibre.  At each of its two ends the
attachment disk of a sufficiently small central prism lies in one of
the fixed exact neighborhoods.  Choose an elongated affine box in
the open face prism: its transverse inner rectangle contains the
central cross-section with positive margin; its longitudinal inner
interval contains the portion outside those two exact neighborhoods;
and both longitudinal end strips of its larger box lie strictly inside
the exact neighborhoods.  The larger box remains strictly between
the two endpoints of the full normal interval.  It therefore avoids
the full blocks themselves.  First decrease $r_{u,D}$, relative also
to the aspect of $D$, to put this whole transverse construction in
$r\ll a_D$; then choose the shell thickness in normalized box
coordinates.  Decrease these transverse and shell sizes also to meet
the fixed $\tau/10$ motion bound: in an affine box the lift displacement
and every hole-buffer diameter are bounded by a fixed box constant
times the normalized shell size, while the inner block will be the
identity.  This bound applies to all the later strict fills as their
images lie in those same buffers.  The two end margins absorb both
that thickness and the sampling enlargements required below.  Distinct boxes and their
sampling neighborhoods are disjoint when they belong to distinct
central prisms, and the neighborhoods form a locally finite family.
This follows by choosing them inside the open face products and
using their compact attachment margins.  The tube portions beyond
the inner longitudinal interval are already exact.

Here and below a sampling neighborhood includes a little more than
the immersion image itself.  The enlargement is needed for the
degree tests in the fibre argument.  All these neighborhoods are
chosen with compact closure in the indicated face and end regions.
Now choose a positive $1$-Lipschitz minorant $\rho\leq\tau/10$ of
all the resulting target value-closeness requirements, including their
affine normalization factors.  Open $\widehat F$ by
Lemma~\ref{hp:opening} and apply Lemma~\ref{sm:relative-pl}, relative
to the fixed enlarged protection balls, to obtain a PL homeomorphism
$h_*$ satisfying
\[
 |h_*(x)-\widehat F(x)|<\rho(\widehat F(x))/10.
\]
Define $K_0=\widehat Fh_*^{-1}$.  At $y=h_*(x)$, the Lipschitz
property of $\rho$ gives
$|K_0(y)-y|<\rho(y)/9$.  Thus all previously fixed sampling and
properness tests hold.  The map $K_0$ is proper, locally Sobolev,
exactly the identity on the protected neighborhoods, and inherits
the old relative opening property.  Its local Sobolev constants may
depend on $h_*$; only their finiteness is used below.

The motion bounds already chosen imply, for every resulting map $g$
and its sufficiently close homeomorphic approximants,
\begin{equation}\label{qt:central-proper-localization}
 g^{-1}(C_m)\subset K_0^{-1}(C_{m+1}).
\end{equation}
Indeed, $|g(y)-K_0(y)|<\tau(y)/5$, whereas
$\tau(K_0(y))>9\tau(y)/10$.  If $K_0(y)\notin C_{m+1}$, the defining
bound for $\tau$ therefore prevents $g(y)$ from lying in $C_m$.
All these requirements preceded the choice of $h_*$.

\emph{The lifted replacement.}
We describe a single normalized box.  Let $B=[-1,1]^3$ be its inner
block, and use the shell, mesh, omitted balls $V_{\sigma}$, buffer balls
$V_{\sigma}^+$, and skeleton immersion $i$ of Lemma~\ref{hp:cutoff}.
Here $\sigma$ indexes shell tetrahedra; $D$ continues to denote a
face subcell.  At depth scale $l$, the immersion has chart radii
comparable to $l$ in its domain
and to $s$ in its image.  On each injective chart $U$,
\[
 |Di|\le C s/l,\qquad |D(i|_U)^{-1}|\le C l/s.
\]
Prescribe $i=\mathrm{id}$ on a whole outer mesh band, leave that
band free of cap supports, and retain the old map there.  On all
tetrahedra whose fills are needed in the protected longitudinal end
strips, put the immersion samplings inside the exact region of
$K_0$.  Their lifts and fills are then the identity on an open
neighborhood, including the necessary collars.  This is a relative
use of the skeleton construction in Lemma~\ref{hp:cutoff}.

On the remaining usable lift region define the near-point branch
\begin{equation}\label{qt:central-lift}
 H(x)=(i|_{U_x})^{-1}\bigl(K_0(i(x))\bigr),
 \qquad i(H(x))=K_0(i(x)).
\end{equation}
The sampling requirements give
$|K_0-\mathrm{id}|\le\varepsilon_{\rm lift}s$, where
$\varepsilon_{\rm lift}>0$ is a fixed closeness constant below all
chart and buffer margins.  Hence
$|H(x)-x|\le C\varepsilon_{\rm lift}l$.
The near-point branches agree on overlaps;
their existence and agreement do not require injectivity of $K_0$.
Neither does the chain-rule estimate in the cutoff proof.  On each
tetrahedron it bounds the lift energy by
\[
 C_p(l/s)^3
       \int_{\text{enlarged image sampling}} |DK_0(y)|^p\,dy.
\]
Summing over depth scales gives the finite shell bound of that
proof.  Affine normalization introduces finite constants depending
on this fixed elongated box, which is harmless here.

For every nonidentity omitted ball, use the elementary forward
squeeze $J_{\sigma}$ which collapses $V_{\sigma}^+$ to its center $v_{\sigma}$, is a
homeomorphism off that closed ball onto the punctured output
region, and is the identity near the tetrahedron boundary.
The lifted boundary of the omitted hole lies strictly inside
$V_{\sigma}^+$.  Define the new map as $J_{\sigma}H$ where the lift is used and
as the constant $v_{\sigma}$ throughout the omitted hole.  These definitions
agree on a whole collar.  Identity-filled holes use neither a cap
nor a nonidentity fill.  Write $J_{\rm cap}$ for the disjoint union
of the squeezes.  Set the map equal to the identity on $B$ and to
$K_0$ beyond the replacement.  This defines $\overline K$.

Continuity at the inner seam follows since all displacements on a
depth-$l$ tetrahedron are $O(l)$.  More precisely, the near-branch
lift and every cap stay on the exterior side of $B$.  Indeed the
Whitney scale comparability gives $\dist(\sigma,B)\geq\alpha l$ for one
fixed $\alpha>0$.  Our choice of $\varepsilon_{\rm lift}$ includes
$C\varepsilon_{\rm lift}l<\alpha l/2$ in the lift bound.
Every lifted point is then strictly exterior to $B$.  Each
cap is supported in an exterior tetrahedron and maps its support into
itself, so cap postcomposition and the omitted-hole constants are
strictly exterior as well, at every depth scale.  The summed lift estimates, followed by the Sobolev gluing argument in
Lemma~\ref{hp:cutoff}, prove local $W^{1,p}$ regularity.  Constant
holes cost no energy.  At the outer band the immersion is the
identity and no cap acts, so the seam is unchanged on a neighborhood.
The reserved outer band and sufficiently small closeness bounds also
keep cap supports and their inverse action inside the replacement
image; hence no source outside the box is changed.  Local finiteness
now proves these assertions simultaneously for all boxes.  Their
low-region margins ensure that both old and new actual labels stay
below $a_D$.

\emph{Homeomorphic openings and connected fibres.}
Apply the old opening property to
obtain homeomorphisms $k_n\to K_0$ finely, preserving the exact
regions and end-strip samplings.  Lift $k_n$ by the same near-point
branches.  The homeomorphic cutoff argument of
Lemma~\ref{hp:cutoff} fills the omitted balls topologically, with
images inside $V_{\sigma}^+$ and the required boundary values.  Replace
each cap by a strictly increasing radial squeeze tending to $J_{\sigma}$.
The resulting maps $\bar k_n$ are homeomorphisms, agree with the
exact data, and preserve the outer replacement image.
On every omitted hole their outputs converge to the declared
constant, irrespective of the choice of topological fill.  Elsewhere
the lift formula gives convergence; at the inner seams the $O(l)$
bound gives uniform convergence by first discarding sufficiently
small depth scales.  Local finiteness gives convergence on compacts,
and diagonal choices give any prescribed fine accuracy.

By \eqref{qt:central-proper-localization}, the inverse images under
all sufficiently close $\bar k_n$ of a fixed compact lie in one
compact.  In particular $\overline K$ is proper and onto: take a
convergent subsequence of $\bar k_n^{-1}(y)$ for a given $y$.
Its fibres are connected as well.  For two points $a,b$ in a fibre
over $y$, choose closed balls $\overline B(y,\rho_n)\Subset\Lambda$,
with $\rho_n\downarrow0$, containing both $\bar k_n(a)$ and
$\bar k_n(b)$.  Their inverse images are compact connected sets
containing $a,b$ in a common compact.  A Hausdorff-convergent
subsequence has connected limit in $\overline K^{-1}(y)$ containing
$a,b$.  This proves the claim.

\emph{Local fibre identification and exact boundary data.}
The following formula gives countability of the new exceptional
values and transfers the wall tests in the next lemma to old
relative-opening tests.  Suppose $y$ is not a cap center and set
$x'=J_{\rm cap}^{-1}(y)$.  In a nonidentity lift region this inverse
is unique, and $x'$ lies outside the closed collapsed buffers
$V_\sigma^+$.  Their fixed clearance from the omitted holes
$V_\sigma$ gives a ball
$B_{\rm s}=B(x',\alpha l)$ in the usable lift region, with an
injective immersion chart on a larger neighborhood $U$.  The joint
immersion neighborhood of Lemma~\ref{hp:immersed-skeleton} supplies
these charts also across ordinary skeleton seams; the limiting inner
seams $S$ are treated separately below.  Choose
$B_{\rm t}=B(i(x'),\beta s)$ inside the enlarged image sampling
neighborhood, reducing the fixed $\beta>0$ so that its chart
pullback lies in $B(x',\alpha l/2)$.  This is possible by the
inverse derivative bound $C l/s$.

The closeness constant was chosen with
$\varepsilon_{\rm lift}<\beta/2$ and
$C\varepsilon_{\rm lift}<\alpha/2$.  On $\partial B_{\rm t}$,
closeness to the identity gives degree one for $K_0$ at $i(x')$
and excludes that value from the boundary image.  Hence the old
fibre meets $B_{\rm t}$ and misses its boundary.  Connectedness
puts the entire old fibre inside $B_{\rm t}$.  In fact every point
of this fibre lies within
$\varepsilon_{\rm lift}s$ of $i(x')$, by the same closeness test there.

Its chart pullback is therefore well inside $B_{\rm s}$.
Equation~\eqref{qt:central-lift} and the near-branch choice identify
this pullback with the local new fibre.  The relative displacement
bound excludes preimages on $\partial B_{\rm s}$.  Connectedness
of the new fibre, applied once more, excludes any other preimage.
Consequently
\begin{equation}\label{qt:central-fibre-formula}
 \overline K^{-1}(y)
   =(i|_U)^{-1}\!\left(K_0^{-1}
          \bigl(i(J_{\rm cap}^{-1}(y))\bigr)\right).
\end{equation}
All margins can be kept uniformly positive after shrinking to a
small target neighborhood $W$ of $y$.  Thus this formula holds for
\emph{every} value in $W$, with one injective chart $U$, and not
merely for its center.  The same argument crosses skeleton seams;
the outer band uses the identity immersion.

Take a countable cover by these charts.  Outside cap centers,
\eqref{qt:central-fibre-formula} shows that a nonsingleton value
must arise from a chart inverse of an old exceptional value.  There
are countably many such values.  Equivalently, a local diffeomorphism
has discrete, hence countable, point preimages in this second-countable
domain.  Identity-filled regions have open identity behavior and
thus singleton fibres by connectedness.  A point of an inner seam
is also isolated in its fibre: on the inner side the map is the
identity, and nearby exterior outputs stay strictly exterior even
after the cap operation.  Its entire fibre is therefore singleton.
Finally, in an unchanged open region a singleton old fibre gives
an isolated point in the new fibre, so a nonsingleton new fibre
meeting that region must have an old exceptional value.  These
cases exhaust the construction and prove countability.

The identity neighborhood of the removed thick cores has singleton
fibres, again because each such point is isolated in its connected
fibre.  It follows that $\overline K^{-1}(M_y)=M_y$, including the
boundary, and gives the asserted exact boundary data after composing
with $h_*$.  The pullbacks of the affine inner cube boundaries form
the stated locally finite family $S$; the end strips make its
nonexact portions disjoint from a neighborhood of $\partial M_x$.
Finally, $T$ is locally Lipschitz, and on $M_y$ every central radius
is positive.  On compact portions only finitely many product and
sector charts occur.  Their radial and perimeter functions have
finite Lipschitz constants there, and near the boundary $F_*$ is
PL.  Sobolev composition, with local Lipschitz extensions of these
coordinate functions, supplies the claimed finite ambient budgets.
\end{proof}

The later wall estimates count intersections along the carrier's mesh
edge paths.  Opening the carrier must preserve these very hits, or the
counts would no longer describe the pulled-back walls.  The next lemma
does this for the sampled wall families, using their whole-wall
avoidance of nonsingleton values.

\begin{lemma}[Exact edge matching when opening the central carrier]
\label{qt:edge-matching}
Use the carrier and exterior of Lemma~\ref{qt:central-carrier}.
Let $\mathcal E$ be a locally finite family of mesh edges in $M_x$,
and let $\mathcal W$ be a locally finite family of sampled standard
walls of Proposition~\ref{qt:exterior}.  Assume:
\begin{enumerate}
\item the entire union $W=\bigcup\mathcal W$ is relatively closed
in $M_y$ and avoids all nonsingleton values of $\overline K$;
\item the actual hit set
$H=(\bigcup\mathcal E)\cap F_*^{-1}(W)$ is finite on compact
localizations, and no hit outside the exact regions lies on $S$.
\end{enumerate}
Then there are homeomorphisms $f_n:\Omega\to\Lambda$ converging
finely to $F_*$, equal to $F_*$ on a neighborhood of $\partial M_x$
and on the other prescribed exact regions, for which
\[
 f_n^{-1}(W_0)\cap e=F_*^{-1}(W_0)\cap e
 \quad\text{for every }W_0\in\mathcal W, e\in\mathcal E.
\]
In particular each restriction maps $M_x$ onto $M_y$ with the exact
prescribed boundary map.  The hypotheses concern sampled edge
tests; they make no relative-opening assertion for arbitrary new
closed tests accumulating through the inner lift scales.
\end{lemma}

\begin{proof}
Work first in the target-domain coordinates given by $h_*$.
Consider a hit $x$ outside an exact region, and write
$y=\overline K(h_*(x))$.  Because $W$ avoids all new nonsingleton
values, $y$ is not a cap center.  The seam-avoidance assumption
permits a target neighborhood $V$ in which the chart formula
\eqref{qt:central-fibre-formula} holds with fixed margins.  For each
wall meeting $y$, take a compact patch $P$ of that wall in $V$,
containing a relative neighborhood of $y$.  The formula shows that
\begin{equation}\label{qt:central-old-wall-test}
 C_P=i\bigl(J_{\rm cap}^{-1}(P)\bigr)
             \quad\text{is compact and satisfies}\quad C_P\cap E_0=\varnothing.
\end{equation}
Indeed, a nonsingleton old fibre at any value of $C_P$ would give
a nonsingleton new fibre at the corresponding point of $P$, which
is excluded by the hypothesis on the \emph{whole wall}.

If the hit lies in an unchanged open region, the same conclusion
uses a direct patch $P$.  Its old fibre at $y$ is singleton:
otherwise the old connected fibre would have an isolated point
where the new fibre is singleton.  Properness then places the old
inverse image of a sufficiently small neighborhood of $y$ entirely
inside the unchanged region.  To see this, a contrary sequence of
values tending to $y$ and preimages outside that region has a
convergent preimage subsequence, contradicting the singleton fibre.
For every value on the smaller wall patch the old fibre consequently
agrees with the new one and is singleton.  Hits in exact open
regions require no additional tests.

There are only finitely many hits on each compact localization.
Thus their neighborhoods can be chosen locally finitely, avoiding
$S$ wherever required.  Only finitely many depth scales occur in
each compact collection of these charts.  Choose the old patch
images in the compactly localized sampling neighborhoods used in
the preceding Lemma.  The resulting family of compact sets $C_P$
is locally finite in $\Lambda$; in particular its union is relatively
closed.  Adjoin the closed protected exact data, slightly inside
their exact neighborhoods.  All these old target tests avoid $E_0$.
Apply Lemma~\ref{hp:opening} to obtain arbitrarily close old
homeomorphisms $k_n$ agreeing with $K_0$ over every one of these
tests and retaining the exact strips.  Here agreement over $C_P$
means the exact old evaluations on $K_0^{-1}(C_P)$; since these
fibres are singleton and $k_n$ is bijective, it also means equality
of the inverse evaluations.

Lift and fill these homeomorphisms as before.  For $y\in P$ the
fixed branch and \eqref{qt:central-old-wall-test} give exactly the
same inverse image of $J_{\rm cap}^{-1}(y)$ under the new lift as
under the carrier lift.  Choose the strict radial cap openings so
that their inverse maps agree with $J_{\rm cap}^{-1}$ on all tested
patches.  This is possible by changing a collapse only over an
arbitrarily small ball about its center: each center has positive
distance from the relatively closed union of the tests relevant
to it, although no uniform distance is required.  The ordinary
piecewise radial interpolation is strictly increasing inside that
ball and unchanged outside it.  Null mesh sizes at the limiting
seams and the reserved outer bands retain the convergence and
support conclusions of the preceding proof.  The resulting
homeomorphism therefore has exactly the carrier inverse over every
tested wall patch; global uniqueness follows from its bijectivity.

Shrink the source hit neighborhoods so their old images, and then
their sufficiently close new images, meet walls only in the chosen
patches.  On the rest of the edges, fine closeness prevents new
hits.  Explicitly, after removing these neighborhoods, the edge
remainder on each compact localization has compact image disjoint
from the relatively closed wall union, hence a positive gap.
An exhaustion and a positive minorant impose all these local gap
conditions simultaneously.  This proves the claimed equality for
every wall and edge, and composition with $h_*$ returns to source
coordinates.  The exact boundary neighborhood is retained throughout,
so the homeomorphisms preserve $M_x$ and its prescribed boundary
map as asserted.
\end{proof}

\begin{remark}
The edge hypotheses can be obtained by the later ACL and coarea
sampling step: choose edges for which the relevant label paths are
absolutely continuous, arrange that their intersections with the
nonexact inner seams have parameter measure zero, and avoid the
resulting null sets of scalar labels when sampling levels.  Finite
variation gives finitely many hits at almost every sampled level.
Countably many exceptional values exclude only countably many
radial or intermediate meridian levels; a graph vertex is not an
intermediate meridian value.  These sampling facts are hypotheses
of the preceding Lemma, rather than an assertion that arbitrary
edges or arbitrary sampled levels satisfy them.
\end{remark}
